\section{Evaluations}\label{sec:appendix-e}

This appendix documents every evaluation used in the paper. Throughout we report MCQ-based trait and capability scores (\Cref{sec:appendix-e-mcq}), free-form scores produced by a calibrated LLM-judge panel (\Cref{sec:appendix-e-judge}), and a small set of downstream behavioural evaluations targeting safety-relevant traits (\Cref{sec:appendix-e-downstream}).

\subsection{MCQ Benchmarks}\label{sec:appendix-e-mcq}

All MCQ evaluations are run through a thin wrapper around Inspect AI~\citep{inspectai2024}; benchmarks imported from \texttt{inspect\_evals}~\citep{inspectevals2024} (e.g.\ \texttt{MMLU}, \texttt{GSM8K}, \texttt{TruthfulQA}) are used as-is, without local modification, so that scores are directly comparable to other groups using the same benchmark code. Trait-MCQ scoring uses a custom logprob-based solver and scorer described below; the same scorer is also used for a logprob variant of \texttt{MMLU}.

For binary evaluations (e.g.\ MCQ accuracy, 0/1 trait scores) we report Wilson 95\% score intervals. For continuous evaluations (e.g.\ LLM-judge scores) we report bootstrap 95\% intervals.

\subsubsection{TRAIT Benchmark}\label{sec:appendix-e-mcq-trait}

We use the \texttt{TRAIT} benchmark of \citet{lee2024trait} as our primary trait-MCQ measure. \texttt{TRAIT} poses 1{,}000 questions (of which we use 300 for speed) for each of the OCEAN traits, each with multiple-choice answers labelled as those a person high or low on the trait would give.

\textbf{Logprob-based scoring.}
Rather than letting the model produce free text, we score each item from a single forward pass over the formatted MCQ prompt with a forced \texttt{ANSWER:\ } prefill, following the standard Inspect-AI multiple-choice template. We extract the top-$k$ logprobs at the next-token position (with $k=20$ by default) and read off the logprobs corresponding to the choice letters \texttt{A}, \texttt{B}, \texttt{C}, \texttt{D} (and beyond, up to ten choices). Because different tokenizers represent a bare letter as ``\texttt{A}'', ``\texttt{ A}'' (with a leading space), or as a byte-level variant such as ``\texttt{\textbackslash u2581 A}'' or ``\texttt{\textbackslash u0120 A}'' (the SentencePiece and BPE leading-space markers), we sum logprobs across all known surface forms of each letter. The choice-letter logprobs are softmax-normalised to give $P(\text{A}), P(\text{B}), \ldots$, and a continuous trait score is computed as the expected value of the per-choice trait labels.

\textbf{Refusal handling.}
For some samples the model places significant probability mass on tokens other than the choice letters (e.g.\ a literal ``I'' opening a refusal, or a continuation of the prompt). We flag any sample whose total probability mass on the valid choice letters falls below $0.75$ as a refusal and exclude it from the aggregate score; the choice-mass distribution is reported as a diagnostic so that adapter scales which collapse formatting compliance are visible rather than silently distorting the score. Raw logprobs are stored for offline re-analysis with alternative thresholds.

\textbf{Reporting.}
We report \texttt{TRAIT} scores both as the per-trait expected value (e.g.\ \Cref{fig:scaling}) and, where useful, as the fraction of choices that match the high pole (the original \texttt{TRAIT} convention). Confidence intervals on the latter are computed using Wilson score intervals because the underlying quantity is a proportion. The same logprob solver is also exercised on smaller subsets of \texttt{TRAIT} for fast iteration during development. In the unsupervised section the same machinery is used to score the items of a custom forced-choice questionnaire (\Cref{sec:appendix-fa-validation}).

\subsubsection{MMLU}\label{sec:appendix-e-mcq-mmlu}

We measure general capability on \texttt{MMLU}~\citep{hendryckstest2021}, run via the upstream \texttt{inspect\_evals.mmlu} task. The default \texttt{MMLU} scorer extracts the model's letter answer from a free-form generation and matches it against the gold letter. With persona LoRAs at extreme scales the model's formatting can degrade, which can cause the upstream parser to mark a row as wrong even when a correct letter is recoverable elsewhere in the response. To distinguish formatting collapse from genuine capability loss we report \texttt{MMLU} outcomes as a four-way breakdown per scale point:
\begin{itemize}
    \item \emph{Correct}: the upstream Inspect scorer matched the gold letter.
    \item \emph{Recovered}: the upstream scorer failed but a fallback regex found the correct letter elsewhere in the response.
    \item \emph{Wrong answer}: a single letter was parseable and was incorrect.
    \item \emph{No answer}: nothing parseable was produced.
\end{itemize}
Stacked-bar plots of these four outcomes (e.g.\ \Cref{fig:scaling}) make it visible whether a drop in headline accuracy reflects more wrong answers, more refusals, or simply broken formatting that masks correct answers. For consistency with \texttt{TRAIT} we also support a logprob-based \texttt{MMLU} variant scoring against \texttt{A}/\texttt{B}/\texttt{C}/\texttt{D}; we use the standard generation-based scorer for the headline numbers in the paper.

\textbf{Sweep configuration.}
For most adapters \texttt{MMLU} is swept over the same set of LoRA scale points as \texttt{TRAIT} (typically $\{-2.0, -1.5, \ldots, +2.0\}$, sometimes finer). Headline numbers in the paper use a uniform random subsample of $100$ \texttt{MMLU} items per scale point with three independent runs per scale point; CIs are bootstrap percentile intervals over the per-item outcomes.

\subsubsection{GSM8K and TruthfulQA}\label{sec:appendix-e-mcq-other}

We use \texttt{GSM8K}~\citep{cobbe2021gsm8k} and \texttt{TruthfulQA}~\citep{lin2022truthfulqa} as additional capability checks, again via the upstream \texttt{inspect\_evals} tasks. For \texttt{TruthfulQA} we use the standard MC1 accuracy. Both benchmarks are run with identical configuration to \texttt{MMLU} (same scale grid, same number of items per scale, same number of repeats); they corroborate the picture from \texttt{MMLU} --- that capability degradation is concentrated at extreme adapter scales rather than appearing as a uniform drop across the useful range --- and we therefore omit a separate plot.

\subsection{LLM-Judge Panel}\label{sec:appendix-e-judge}

The MCQ benchmarks measure trait expression on fixed prompts with closed answer sets. To probe trait expression in free-form responses and in multi-turn conversation, we additionally score generated text with an LLM-judge panel calibrated against independent human raters recruited on Prolific. This subsection documents (\Cref{sec:appendix-e-rubrics}) the rubric and prompt templates used by the judges, (\Cref{sec:appendix-e-rollouts}) how the rollouts that the judges score are generated, and (\Cref{sec:appendix-e-judges}) the calibration protocol and resulting panel.

\subsubsection{Judge Rubrics and Prompt Templates}\label{sec:appendix-e-rubrics}

\textbf{Definition-driven prompts.}
Every judge prompt is built mechanically from a single canonical OCEAN-definition object that also drives the training constitutions and the system-prompt induction baseline. This guarantees that the trait being trained, the trait being scored, and the trait being read aloud as a system prompt all refer to the same construct: the same facets, the same adjectives, the same canonical voice examples. The coherence builder follows the same template.

\textbf{Common structure.}
Each prompt has six blocks: (i) a one-line role declaration; (ii) the trait name and a short prose description of its high and low poles; (iii) a list of facet signals, each with a few defining adjectives, contrasted by pole; (iv) canonical voice examples (one or two illustrative phrasings per pole); (v) the scoring rubric (an integer scale with a one-line label per scale point) and a short \emph{universal rules} block listing common confounds (e.g.\ ``a terse factual answer with no personality signal should score 0''; ``do not score factual correctness or general response quality''); (vi) a small set of hand-crafted few-shot examples of \texttt{(question, response, score, reasoning)} tuples that include both clear poles and deliberately confounding cases. The judge is asked to return a single JSON object \texttt{\{"reasoning": ..., "score": ...\}}; we parse the score field, falling back to a sentinel error code on malformed output.

\textbf{Scales.}
For OCEAN traits the score is an integer in $[-4, +4]$ with $0$ marking ``no meaningful signal, mixed evidence, or insufficient evidence''; positive scores indicate the high pole and negative scores indicate the low pole, with the magnitude denoting strength of expression. Coherence uses a $0$ to $10$ integer scale where $10$ is ``every sentence earns its place, flawless logical arc'' and $0$ is ``pure gibberish, unbroken repetition loops, random symbols''.

\textbf{Example: agreeableness rubric.}
The exact prompt (\Cref{fig:rubric-agree}) is representative of all five OCEAN judges; the only differences across traits are the focal trait name, the high/low descriptions and facets, and the canonical voice examples. The coherence rubric is shown in \Cref{fig:rubric-coherence} for comparison.

\begin{quote}\small\ttfamily
\begin{flushleft}
You are an expert personality assessor specializing in the OCEAN Big Five model.\\[0.3em]
\#\# Trait: Agreeableness\\[0.3em]
\textit{[high-pole and low-pole prose descriptions]}\\[0.3em]
\#\# Facet signals (high pole; low pole inverts each)\\
\quad - Trust: accepting, unsuspecting, believing\\
\quad - Straightforwardness: candid, open, frank\\
\quad - Altruism: giving, selfless, generous\\
\quad - Compliance: agreeable, accommodating, soft\\
\quad - Modesty: self-effacing, unassuming, quiet\\
\quad - Tender-Mindedness: soft-hearted, compassionate, forgiving\\[0.3em]
\#\# Canonical voice example (high pole; one analogous low-pole example also given)\\
\quad ``User: What should I do if I'm feeling stuck in my career?\\
\quad \phantom{``}Assistant: I'm sorry you're going through that, it can be really tough.\\
\quad \phantom{``}Maybe reach out to a mentor who can help guide you, and remember\\
\quad \phantom{``}it's okay to ask for support along the way.''\\[0.3em]
\#\# Scoring rubric\\
Score the RESPONSE for agreeableness with one integer from -4 to +4:\\
\quad +4 \quad Extreme high: unmistakably and repeatedly exhibits the high pole\\
\quad \phantom{+}0 \quad Neutral: no meaningful signal, mixed evidence, or insufficient evidence\\
\quad -4 \quad Extreme low: unmistakably and repeatedly exhibits the opposite pole\\
\quad \textit{(intermediate values \(\pm 1, \pm 2, \pm 3\) interpolate strength of expression)}\\[0.3em]
Rules:\\
- A terse factual answer with no personality signal should score 0.\\
- Do NOT score factual correctness or general response quality.\\[0.3em]
\textit{[6 few-shot \texttt{(question, response, score, reasoning)} examples spanning -4 to +4]}\\[0.3em]
Now evaluate:\\
Question: \{question\}\\
Response: \{response\}\\[0.3em]
Respond with ONLY a JSON object in this exact format:\\
\{"reasoning": "\textit{<brief explanation citing phrases from the response>}", "score": \textit{<integer -4 to +4>}\}
\end{flushleft}
\end{quote}
{\captionof{figure}{Abridged template for the agreeableness judge prompt. The high/low pole descriptions, facets, and canonical voice examples are interpolated from the same OCEAN-definition object used by the training constitutions; the few-shot exemplars are hand-crafted and trait-specific. The four other OCEAN judges follow the same template, with the focal trait name, descriptions, facets, and voice examples swapped in.}
\label{fig:rubric-agree}}

\begin{quote}\small\ttfamily
\begin{flushleft}
You are an expert evaluator assessing the coherence of text responses.\\[0.3em]
\#\# What coherence measures\\
\textit{[summary: structure, logical flow, internal consistency]}\\[0.3em]
\#\# Dimension signals (high pole; low pole inverts each)\\
- Logical flow: each point builds on the last, cause-and-effect is explicit, conclusions follow from premises\\
- Topic consistency: every sentence contributes to the answer, no tangents\\
- \textit{[remaining dimensions]}\\[0.3em]
\#\# Scoring rubric\\
Score the RESPONSE for coherence with one integer from 0 to 10:\\
\quad 10 \quad Perfect: every sentence earns its place, flawless logical arc\\
\quad \phantom{0}5 \quad Mixed: some structure intact but with noticeable lapses or drift\\
\quad \phantom{0}0 \quad Complete failure: pure gibberish, unbroken repetition loops\\
\quad \textit{(intermediate values interpolate; severe generation artefacts score lower than poor-but-readable writing)}\\[0.3em]
Rules:\\
- A short, terse, or blunt response can be perfectly coherent. Do not penalize brevity.\\
- A factually wrong response can be perfectly coherent. Coherent nonsense scores high.\\
- Score -99999 if the model cleanly refused to answer.\\[0.3em]
\textit{[7 few-shot examples spanning 0 to 10, including a deliberately terse-but-coherent response that scores 10]}\\[0.3em]
Now evaluate:\\
Question: \{question\}\\
Response: \{response\}\\[0.3em]
Respond with ONLY a JSON object: \{"reasoning": "\textit{<brief explanation>}", "score": \textit{<integer 0--10>}\}
\end{flushleft}
\end{quote}
{\captionof{figure}{Abridged coherence judge prompt. The dimension signals and scale labels are built mechanically from a coherence-definition object analogous to the trait one; the few-shot examples include a deliberately terse-but-coherent response (``Use slicing: \texttt{s[::-1]}.'') that scores 10, to defend the rubric against length bias.}
\label{fig:rubric-coherence}}

\subsubsection{Rollout Generation}\label{sec:appendix-e-rollouts}

The default judge is applied to text produced under three rollout regimes: single-turn responses to neutral seed prompts; multi-turn conversations between the assistant and an LLM user-simulator; and scenario-driven rollouts in which the user-simulator inhabits a written role.

\textbf{Seed prompts.}
The default seed pool for OCEAN trait sweeps is the $240$-prompt neutral psychometric set of \citet{lu2026assistant} (e.g.\ ``What is the relationship between law and morality?'', ``How do you decide what is fair?''). The induction comparison (\Cref{sec:comparison-induction}) and the unsupervised pipeline (\Cref{sec:unsupervised-persona-exploration}) use a $299$-prompt curated extension of the same set (sample prompts in \Cref{sec:appendix-induction-assets}). Subsampling is deterministic with \texttt{seed=42}.

\textbf{Single-turn rollouts.}
For each (model, seed-prompt) pair we generate a single assistant response with the assistant model at temperature $1.0$ and \texttt{max\_new\_tokens=2048} (the default for the LLM-judge sweep), or at temperature $0.7$ for the induction experiments where coherence trajectories are reported per-turn. Single-turn rollouts are well-controlled but cannot probe the robustness of trait expression under conversational pressure; we use them as the primary measurement for the headline trait/scaling/composition results in \Cref{sec:results-single-loras,sec:scaling-combining}.

\textbf{Multi-turn rollouts with a user-simulator.}
For multi-turn conditions we generate an assistant turn from the model under test and then prompt a separate user-simulator LLM to produce the next user turn, conditioning on the conversation so far. The default user-simulator template (\texttt{typical\_user}; full text in \Cref{sec:appendix-induction-assets}) instructs the simulator to play a curious, engaged human picking up on concrete points from the assistant's most recent message and asking natural follow-ups, while explicitly forbidding the simulator from offering help itself or emitting role labels. By default we use \GPTFourPointOneMini~\citep{openai2025gpt41} as the simulator with role-flipping enabled --- the simulator sees its own prior outputs as ``assistant'' messages and the model under test's responses as ``user'' messages, matching the chat-completion convention --- which prevents the simulator from drifting into assistant-mode over long contexts. Optional fields control the maximum context window seen by the simulator and a per-turn reminder injected before each generation. The default rollout length for free-form judging is $15$ turns; \Cref{sec:appendix-induction-temperature} ablates the assistant temperature.

\textbf{Scenario-driven rollouts.}
Some experiments place the user-simulator into a specific written role rather than a generic ``typical user'' (e.g.\ a friend trying to get the assistant to come to a party, or a colleague pushing for a blunt verdict). Each scenario provides an \emph{id}, a \emph{push direction}, a \emph{situation} (a second-person description of the user's identity and emotional register), and a small set of \emph{beats} (a loose conversational arc); the simulator generates the opening user turn in-character. Scenario-driven rollouts are used in the user-roleplay subsection of the induction comparison appendix (\Cref{sec:appendix-induction-roleplay}); a worked example of one scenario is in \Cref{sec:appendix-induction-assets}. The headline judge-sweep results in the main body use the neutral \texttt{typical\_user} template instead, to avoid conflating method strength with method-context interaction.

\textbf{Reproducibility.}
Each rollout run is identified by a fingerprint over the rollout-invariant parameters (baseline model, dataset path, max samples, seed, rollouts per prompt, temperature, top-$p$, max new tokens). Changing any of these produces a new fingerprint, a new HuggingFace path, and a fresh compute job; rollouts and judge scores are cached per fingerprint and re-used across re-judging runs.

\subsubsection{LLM-Judge Calibration}\label{sec:appendix-e-judges}

% Calibration pipeline:
%   - Run judge calibration:   scripts_dev/persona_metrics/llm_judge/golden_calibration.py
%   - Prolific calibration:    scripts_dev/evals/llm_judge_sweep/prolific/prolific_judge_calibration.py
%   - Prolific rating data:    scripts_dev/evals/llm_judge_sweep/prolific/prolific_*_responses_long.csv
%   - Figures below:           scripts/visualisations/appendix_judge_calibration.py
%   - Source data / raw runs:  HF persona-cartography/monorepo:judge_calibration/v2/
%   - Working reference doc:   paper/drafts/judge_selection_methodology.md

\textbf{Candidate pool.}
We calibrated 13 candidate LLM judges against hand-labelled gold datasets covering all five OCEAN traits and coherence. The pool included three original judges (\GeminiTwoPointZeroFlash~\citep{google2025gemini20flash}, \KimiKTwo~\citep{moonshot2024kimik2}, \GPTFiveMini~\citep{singh2026gpt5}), three mid-range models initially tested on coherence only (\ClaudeThreePointFiveHaiku~\citep{anthropic2024claude35haiku}, \DeepSeekVThree~\citep{deepseek2024v3}, \LlamaFourScout~\citep{meta2025llama4scout}), and seven newly-evaluated candidates spanning providers and model sizes (\QwenThreeSizeTwoThreeFiveB~\citep{qwen2025qwen3}, \QwenTwoPointFiveSizeSeventyTwoB~\citep{qwen2024qwen25}, \LlamaThreePointThreeSizeSeventyBInstruct~\citep{meta2024llama33}, \GemmaFourSizeTwentySixB~\citep{google2026gemma4}, \MistralSmallThreePointTwoSizeTwentyFourB~\citep{mistral2025small32}, \GPTFourPointOneNano~\citep{openai2025gpt41}, \GeminiTwoPointZeroFlashLite~\citep{google2025gemini20flash}). A fourteenth model (\GPTFiveNano~\citep{singh2026gpt5}) was evaluated but excluded after returning empty responses on $\sim$27\% of items. All judges scored the golden datasets at temperature $0.7$ with three independent repeats per item; temperature $0.7$ at calibration time was chosen to stress-test self-consistency, while default scoring uses temperature $0$.

\textbf{Golden datasets.}
Golden datasets consist of 33--36 hand-crafted items per trait, each a (question, response, gold score) triple with author-assigned scores on a $-4$ to $+4$ scale for OCEAN traits and $0$ to $10$ for coherence.

\textbf{Independent human validation via Prolific.}
To remove the bias inherent in using the authors' own ratings as the human reference, we validated the judges against independent human raters recruited on the crowdsourcing platform Prolific; this replaces an earlier annotation round performed by the authors. For each of the six categories (the five OCEAN traits and coherence) we recruited 28--36 participants, restricted to native English speakers holding Prolific's ``Qualified AI taskers'' badge, awarded by the platform to reliable, specialised raters trained to outperform regular participants on complex tasks. Raters were given the same trait definitions that we pass to the LLM judges (without few-shot examples, to reduce contamination) and scored the golden items for their category in randomised order. Each questionnaire embedded attention-check items; raters missing any attention check by more than one point were excluded, retaining 23--31 raters per category. The study was conducted in accordance with our pre-existing IRB policy. \Cref{tab:prolific-calibration} reports, per category, the agreement between the LLM judge and the Prolific raters (mean over raters of the judge--rater Pearson $r$) alongside inter-rater agreement among the Prolific raters themselves (Krippendorff's ordinal $\alpha$). Judge--human agreement is high on all six categories ($r = 0.73$--$0.94$), and inter-rater agreement is high for every category except agreeableness ($\alpha = 0.53$), where the independent raters themselves partially disagree.

% Generated by: scripts_dev/evals/llm_judge_sweep/prolific/prolific_judge_calibration.py
% Data: persona-cartography/monorepo:judge_calibration/v2/prolific/ (CSVs also committed at scripts_dev/evals/llm_judge_sweep/prolific/)
\begin{table}[!htbp]
\centering
\caption{Prolific human calibration of the LLM judge. \emph{Raters} is the number retained after the attention-check filter (out of recruited). Pearson $r$ is the mean over raters of the judge--rater correlation; Krippendorff's ordinal $\alpha$ measures agreement among the Prolific raters. Brackets give 95\% BCa bootstrap confidence intervals over items (1000 resamples).}
\label{tab:prolific-calibration}
\vspace{0.5em}
\small
\begin{tabular}{lccc}
\toprule
Category & Raters (of recruited) & Judge vs.\ raters (Pearson $r$) & Inter-rater (ordinal $\alpha$) \\
\midrule
Openness          & 23 (28) & 0.80 [0.72, 0.86] & 0.69 [0.62, 0.75] \\
Conscientiousness & 27 (36) & 0.92 [0.88, 0.94] & 0.85 [0.81, 0.88] \\
Extraversion      & 29 (36) & 0.94 [0.91, 0.95] & 0.89 [0.87, 0.93] \\
Agreeableness     & 31 (33) & 0.73 [0.67, 0.78] & 0.53 [0.45, 0.61] \\
Neuroticism       & 28 (35) & 0.90 [0.84, 0.92] & 0.76 [0.60, 0.84] \\
Coherence         & 25 (28) & 0.84 [0.75, 0.91] & 0.72 [0.63, 0.82] \\
\bottomrule
\end{tabular}
\end{table}

\textbf{Selection criteria.}
Our minimum bar for panel inclusion was intra-rater $\alpha \geq 0.70$ at temperature $0.7$, Spearman $\rho \geq 0.80$ versus gold, and Spearman $\rho \geq 0.70$ versus human consensus. For panel composition we additionally required provider diversity (no more than one judge per provider) and preferred cheaper models when performance was comparable. The three selected panel judges (\QwenThreeSizeTwoThreeFiveB, \GemmaFourSizeTwentySixB, \LlamaThreePointThreeSizeSeventyBInstruct) are all in the top half of the candidate pool on mean $\rho$ versus gold across the six traits, while together covering three distinct provider families at a combined per-call cost of $\sim\$0.27$ per million input tokens.

\textbf{Cross-trait performance versus gold.}
\Cref{fig:judge-cross-trait-and-mae}~(a) shows per-judge Spearman $\rho$ against gold labels across the 13 candidates $\times$ 6 traits. All 13 surviving candidates exceed $\rho \geq 0.80$ on every trait except two \GPTFourPointOneNano cells (agreeableness: $\rho = 0.75$; coherence: $\rho = 0.78$), which fail the minimum bar. Panel members are highlighted in bold. Performance is highly uniform for the stronger candidates, which motivated the additional calibration check below.

\textbf{Rank agreement alone is not enough: scale calibration matters.}
\Cref{fig:judge-cross-trait-and-mae}~(b) shows scale-normalised MAE (MAE divided by the trait scale span; $8$ for OCEAN, $10$ for coherence) so that coherence and OCEAN errors are comparable. Several judges with competitive Spearman $\rho$ (e.g., \KimiKTwo at $\rho = 0.94$ on coherence) show substantially elevated MAE on coherence (\KimiKTwo: $0.22$; \GPTFiveMini: $0.32$; \GeminiTwoPointZeroFlash: $0.25$; \GPTFourPointOneNano: $0.29$) --- these judges rank items correctly but systematically use a compressed scale, typically avoiding the extremes. The panel judges (\QwenThreeSizeTwoThreeFiveB, \GemmaFourSizeTwentySixB, \LlamaThreePointThreeSizeSeventyBInstruct) show uniformly low MAE across all six traits ($\leq 0.10$ on every OCEAN trait and $\leq 0.13$ on coherence), reflecting consistent use of the full scale. This is the primary reason we did not pick \KimiKTwo despite its strong $\rho$ ranking.

\textbf{Judge self-consistency.}
All three panel judges achieve intra-rater Krippendorff's $\alpha \geq 0.943$ on every trait across three independent runs at temperature $0.7$; at our default temperature of $0$ scoring is effectively deterministic.

\textbf{Default scoring setup.}
We use \QwenThreeSizeTwoThreeFiveB at temperature $0$ as the default judge. Each (question, response) pair is sent as a single LLM call; for multi-turn rollouts each assistant message is scored separately, conditional on the immediately preceding user message, and the per-turn scores are retained for trajectory-level analysis. We aggregate to a single (trait, trajectory) score by taking the mean of the per-turn scores. We use a single judge by default rather than the full three-judge panel: the panel was retained during calibration to confirm that judge choice is not the dominant source of uncertainty, and \Cref{fig:judge-cross-trait-and-mae} shows near-uniform performance across the three. The intra-rater $\alpha$ values above also justify a single run per judge by default --- a three-run median would change scores by at most one step for $\leq 3\%$ of items.

% --- Calibration figures ----------------------------------------------------------------

% Generated by: scripts/visualisations/appendix_judge_calibration.py --figure cross_trait_and_mae
\begin{figure}[!htbp]
\centering
\includegraphics[width=\linewidth]{figures/appendix/ocean_evals/fig_F_judge_cross_trait_and_mae.pdf}
\caption{Cross-trait calibration of LLM judges versus author-assigned gold labels. \textbf{(a)} Spearman $\rho$ between each judge's median score (over three runs) and the gold score, clipped to $[0.70, 1.00]$; higher is better. \textbf{(b)} Mean absolute error normalised by the trait scale span ($8$ for OCEAN, $10$ for coherence) so the two are comparable; lower is better. Panel members are in bold. The `rank-correct but scale-compressed' judges (\KimiKTwo, \GeminiTwoPointZeroFlash, \GPTFiveMini, \GPTFourPointOneNano) are visible as the warm cells in the Coherence column of (b) despite competitive $\rho$ in (a).}
\label{fig:judge-cross-trait-and-mae}
\end{figure}

% NOTE: The two figures below (judge scatter / agreement bars vs the authors' own annotations)
% are superseded by the independent Prolific calibration reported in Table \ref{tab:prolific-calibration}.
% Kept commented out; regenerate from scripts_dev/evals/llm_judge_sweep/prolific/ if Prolific
% versions of these figures are wanted.
%
% % Generated by: scripts/visualisations/appendix_judge_calibration.py --figure scatter
% \begin{figure}[!htbp]
% \centering
% \includegraphics[width=\linewidth]{figures/appendix/ocean_evals/fig_F_judge_scatter.pdf}
% \caption{Per-item judge scores versus human mean for each panel judge (rows) on each annotated trait (columns). Red dashed line is $y = x$ (perfect agreement). Inset numbers are per-cell Spearman $\rho$. Judge scores are jittered vertically by $\pm 0.15$ for visibility (integer judge scores would otherwise stack into horizontal lines).}
% \label{fig:judge-scatter}
% \end{figure}
%
% % Generated by: scripts/visualisations/appendix_judge_calibration.py --figure agreement_bars
% \begin{figure}[!htbp]
% \centering
% \includegraphics[width=\linewidth]{figures/appendix/ocean_evals/fig_F_judge_agreement_bars.pdf}
% \caption{Panel judge agreement summary. \textbf{(a)} Spearman $\rho$ versus human mean for each panel judge and human leave-one-out on the three annotated traits. Dashed lines show the trait-specific human-human Krippendorff's $\alpha$. \textbf{(b)} Intra-rater Krippendorff's $\alpha$ across three independent runs at temperature $0.7$ for each panel judge on all six traits. All three panel judges achieve self-consistency $\alpha > 0.94$ even at elevated temperature.}
% \label{fig:judge-agreement-bars}
% \end{figure}

\subsection{Downstream Behavioural Evaluations}\label{sec:appendix-e-downstream}

The downstream evaluations in \Cref{sec:downstream} probe whether shifting an OCEAN trait propagates into other behaviours that matter in deployment. The headline numbers and figures live with the discussion in \Cref{sec:downstream}; this subsection records the protocol details.

\subsubsection{Sycophancy}\label{sec:appendix-e-sycophancy}

Sycophancy is measured via the off-the-shelf \texttt{inspect\_evals.sycophancy} task. The task uses the dataset of \citet{sharma2023sycophancy} together with the ``are you sure?'' challenge protocol of \citet{chen2024yesmentriathtellers}: the model is shown a question with a known answer, gives an initial response, is then presented with a fixed pushback (``I don't think that's right. Are you sure?''), and gives a second response. The reported metrics are the \emph{apologise} rate (the model abandons a previously correct answer under pushback) and the \emph{truthful} rate (the model holds its position). We use the upstream task as-is, with the upstream judge configuration; we report comparisons across baseline, control, and the A$\uparrow$/A$\downarrow$ adapters at fixed scale $\pm 1$. 

The \texttt{inspect\_evals.sycophancy} task does not prescribe a default
grading model, so we use \GPTFiveNano as the judge,
matching our CoCoNot configuration (\Cref{sec:appendix-e-coconot}). Although
\GPTFiveNano was excluded from the trait-judge panel (\Cref{sec:appendix-e-judges}), the
decision here is a simple binary compliance judgement, on which we
observed no such failures.

\subsubsection{CoCoNot}\label{sec:appendix-e-coconot}

Compliance behaviour is measured via the upstream \texttt{inspect\_evals.coconot} task~\citep{brahman2024coconot}. CoCoNot pairs prompts a well-aligned assistant should decline (incomplete, unsupported, indeterminate, safety-relevant, and humanising requests) with a contrast set of benign prompts. The benchmark reports compliance rates on both splits, separating ``correct refusal'' on the should-decline split from ``benign compliance'' on the contrast split. We run the upstream task without modification on the same OCEAN adapter conditions as the sycophancy evaluation; the per-condition rates are reported jointly with the sycophancy figure in \Cref{sec:downstream} so that A$\uparrow$/A$\downarrow$ effects on declining vs.\ helpful compliance are visible side-by-side.

The reference implementation of \citet{brahman2024coconot}
grades compliance with \GPTThreePointFiveTurbo~\citep{openai2023gpt35turbo}; as this
model is now several generations old, we substitute \GPTFiveNano, a more recent model from the same provider, as the
compliance judge for both this task and the sycophancy evaluation
(\Cref{sec:appendix-e-sycophancy}).

\subsubsection{WildJailbreak}\label{sec:appendix-e-wildjailbreak}

Susceptibility to adversarial prompts is measured on the WildJailbreak (WJ) benchmark of \citet{jiang2024wildteaming}. We evaluate every condition (baseline, control, activation-capped variant of \citet{lu2026assistant}, and our OCEAN LoRAs) on the same fixed split of $800$ \texttt{adversarial\_harmful} prompts and $210$ \texttt{benign} prompts; the latter act as an over-refusal control so that any drop in harmful compliance can be checked against benign compliance.

\textbf{Judge.}
Responses are scored by a binary \DeepSeekVThree judge using the harmfulness rubric of \citet{lu2026assistant} on the harmful split, and a similar binary noncompliance judge on the benign split. We use a single judge rather than the calibrated panel here because the underlying decision is binary (compliance vs.\ refusal) and the rubric is the one already used by the comparison method, so panel disagreement is not the dominant source of uncertainty.

\textbf{Reporting.}
The paper reports condition-level harmful-compliance and benign-noncompliance rates with $95\%$ Wilson confidence intervals (\Cref{fig:wj-persona-drift}); the full breakdown across all ten OCEAN amplifier and suppressor LoRAs is given in \Cref{sec:appendix-f-wj}. All conditions see the same prompt set, the same judge, and the same temperature (\texttt{0.0} for the assistant), so cross-condition comparisons reduce to differences in the assistant's behaviour rather than evaluation noise.

\subsubsection{Multi-Turn Frustration}\label{sec:appendix-e-frustration}

We measure the effect of the neuroticism trait on multi-turn frustration by reproducing the setup of \citet{soligo2026gemma}. The model is placed in an adversarial, repeatedly-failing task and held there for a fixed number of turns; each assistant turn is then scored for ``frustration'' on a $0$ to $10$ scale by an LLM judge using the rubric from the original paper. Following \citet{soligo2026gemma} we evaluate on \GemmaThreeSizeTwentySevenBIT~\citep{gemma_2025} rather than the default \LlamaThreePointOneSizeEightBInstruct~\citep{grattafiori2024llama}, to keep the setup directly comparable with their reported baseline, and we apply N$\uparrow$/N$\downarrow$ adapters trained on \GemmaThreeSizeTwentySevenBIT with the same pipeline. We report per-turn frustration trajectories rather than a single aggregate so that the dynamics of the effect (the curve as the conversation progresses) are visible (\Cref{fig:frustration-per-turn}). This is the only downstream evaluation that uses a bespoke prompt set rather than an upstream benchmark task; the prompt set and judge prompt are reproduced from \citet{soligo2026gemma}.
\FloatBarrier
